\section{Numerical level sets and hidden constraints}

For nonlinear saturation maps, measured flux data, or \gls{fe} results, the
surfaces are no longer exact quadrics.  The analytic model remains a regression
case for numerical geometry:
\begin{enumerate}
 \item sample voltage, torque, and loss on a grid or irregular point cloud;
 \item use triangles in 2D or tetrahedra/voxels in 3D to locate sign changes of
       a requested level;
 \item reconstruct ordered contours or triangulated isosurfaces rather than
       treating edge intersections as an unstructured cloud;
 \item verify physical residuals and gradient/KKT residuals at candidate
       optima.
\end{enumerate}
Delaunay simplices are useful for irregular data; marching-cubes-style
isosurfaces are preferable for clean visualization on regular grids.  These
methods reveal deformed surfaces without requiring an algebraic solution for
one current coordinate.

\subsection{Hidden numerical constraints}

A finite search box silently adds current limits.  In particular, solving
\cref{eq:naive-mtpv} on
\begin{equation}
 i_d\in[i_{d,\min},i_{d,\max}],\quad
 i_q\in[i_{q,\min},i_{q,\max}],\quad
 i_e\in[0,i_{e,\max}]
\end{equation}
can return a finite point only because a box face truncates the voltage
cylinder.  A physical optimum should converge when the numerical domain is
enlarged; an artifact moves with the boundary.  Report which constraints are
active, repeat the solve on expanded domains, and inspect whether the result is
on a mesh boundary.

\paragraph{Small domain-enlargement experiment.}
Solve first on $\abs{i_d},\abs{i_q},\abs{i_e}\leq1$, then repeat with bounds
$2$ and $4$ in normalized units.  A physical optimum should settle while the
box recedes.  If the maximizing point and torque keep moving outward, the
experiment has diagnosed an unbounded model rather than a better operating
point.

For visualization, fixed-excitation slices often communicate the geometry
better than many transparent 3D surfaces.  A few labeled slices, the null
direction, and the loss ellipsoid make compactness and active constraints
visible without obscuring them.
